\chapter{Conclusion}
\label{chap:conclusion}

This dissertation aimed to go beyond just measuring accuracy when
evaluating mathematical reasoning in Large Language Models. It explored
what types of errors three models make, whether longer reasoning leads
to more correct answers, and whether Tree-of-Thought prompting actually
improves results compared to Chain-of-Thought. The answers to these
questions are not straightforward. Gemini, GPT-OSS-120B, and Mistral
Large all share the main weakness of Combinatorial Counting Error,
suggesting this is a common limitation in current mathematical reasoning
methods, not just a problem with one provider. For two of the three
models, longer reasoning does not mean more reliable answers. For
Mistral Large, the exception, incorrect answers actually involve more
reasoning steps than correct ones -- extra effort does not translate
into a better outcome, which may be because the model cannot control how
long its reasoning is. Tree-of-Thought prompting improves accuracy for
two of the three models, but this comes with a considerable
computational cost. For the third model, the same effort does not lead
to any clear benefit.

Overall, these results show that LLM mathematical reasoning should not
be seen as a single, uniform skill that reacts the same way to every
prompting strategy. For people considering Tree-of-Thought prompting in
a mathematical reasoning task, the data show that its effectiveness
depends on the specific model and should be tested directly, not assumed
based on results from other models. It is important to evaluate each new
model or prompting method on its own, instead of expecting patterns from
previous studies always to apply.

This study took place at a specific time, using three free-tier models
tested with a curated set of 131 questions. Some parts of the analysis
have real constraints: the neutral-judge validation was left unfinished
due to the provider's daily token quota, and the RQ2 significance test,
while run as a supplementary check, gave mixed, test-dependent results
rather than a confirmed finding, as explained in
Section~\ref{sec:threats-validity}. Even with these limits, the main
conclusion remains: looking only at
accuracy would have missed the shared error pattern, the link between
step count and correctness for each model, and the model-specific
effects of Tree-of-Thought prompting that this deeper analysis showed.
